\chapter{تشخیص داده‌های پرت}
\label{ch:outliers}

\section*{اهداف این فصل}
\begin{itemize}
    \item درک مفهوم داده پرت و تفاوت آن با نویز
    \item آشنایی با روش‌های آماری ساده برای تشخیص داده پرت
    \item آشنایی با روش‌های مبتنی بر فاصله و چگالی
    \item آشنایی با روش‌های مبتنی بر یادگیری ماشین
\end{itemize}

\section{تعریف و انواع داده‌های پرت}

همان‌طور که در فصل \ref{ch:data} اشاره شد،
\dmterm{\textbf{داده پرت}}{Outlier} نمونه‌ای است که به‌طور قابل‌توجهی با
بقیهٔ داده تفاوت دارد. برخلاف نویز که معمولاً خطای بی‌معنی و قابل‌حذف
است، داده پرت می‌تواند خودِ هدف اصلی تحلیل باشد: تراکنش متقلبانه در
داده‌های بانکی، رفتار غیرعادی یک حسگر پیش از خرابی تجهیزات صنعتی، یا
یک بیمار با علائم نادر و نگران‌کننده. به همین دلیل، این وظیفه گاهی
\dmterm{\textbf{تشخیص ناهنجاری}}{Anomaly Detection} نیز نامیده می‌شود.

داده‌های پرت را می‌توان به سه دسته تقسیم کرد:

\begin{itemize}
    \item \dmterm{\textbf{پرت سراسری}}{Global Outlier}: نمونه‌ای که در
          مقایسه با \textbf{کل} مجموعه داده غیرعادی است (مثلاً دمای بدن
          ۴۲ درجه).
    \item \dmterm{\textbf{پرت زمینه‌ای}}{Contextual Outlier}: نمونه‌ای
          که تنها در یک زمینهٔ خاص غیرعادی است (مثلاً دمای هوای ۳۵ درجه
          در زمستان غیرعادی است، اما در تابستان طبیعی است).
    \item \dmterm{\textbf{پرت جمعی}}{Collective Outlier}: مجموعه‌ای از
          نمونه‌ها که هرکدام به‌تنهایی عادی به نظر می‌رسند، اما رفتار
          جمعی آن‌ها غیرعادی است (مثلاً یک الگوی ترافیک شبکه که نشانهٔ
          حملهٔ سایبری است).
\end{itemize}

\section{روش‌های آماری}

ساده‌ترین رویکرد برای تشخیص داده پرت، استفاده از فرضیات آماری دربارهٔ
توزیع داده است. برای داده‌ای که تقریباً توزیع نرمال دارد، رایج‌ترین
روش استفاده از \dmterm{\textbf{امتیاز} \LR{Z}}{Z-Score} است (که پیش‌تر
در فصل \ref{ch:data} برای استانداردسازی دیدیم): نمونه‌ای که قدرمطلق
امتیاز $Z$ آن بیش از یک آستانه (معمولاً ۲.۵ یا ۳) باشد، به‌عنوان داده
پرت علامت‌گذاری می‌شود.

روش دیگر که وابستگی کمتری به فرض نرمال بودن توزیع دارد، استفاده از
\dmterm{\textbf{دامنهٔ میان‌چارکی}}{Interquartile Range (IQR)} است: هر
مقداری که خارج از بازهٔ
$[Q_1 - 1.5 \times IQR,\ Q_3 + 1.5 \times IQR]$ قرار گیرد (که در آن
$Q_1$ و $Q_3$ به‌ترتیب چارک اول و سوم داده‌اند)، به‌عنوان پرت در نظر
گرفته می‌شود. این روش پایهٔ نمودار جعبه‌ای است\footnote{\LR{Box Plot}}.

محدودیت اصلی روش‌های آماری این است که معمولاً فقط برای یک ویژگی در یک
زمان مناسب‌اند (تک‌متغیره) و به‌سختی به داده‌های چندبعدی با روابط
پیچیده میان ویژگی‌ها تعمیم می‌یابند.

\section{روش‌های مبتنی بر فاصله و چگالی}

روش‌های مبتنی بر فاصله، داده پرت را نمونه‌ای تعریف می‌کنند که همسایگان
نزدیک کمی دارد یا فاصلهٔ آن تا نزدیک‌ترین همسایه‌هایش زیاد است — رویکردی
مشابه چیزی که در DBSCAN (فصل \ref{ch:clustering}) دیدیم، که در آن نقاط
کم‌چگالی به‌طور طبیعی به‌عنوان نویز شناسایی می‌شدند.

روش دقیق‌تری که چگالی \textbf{نسبی} را در نظر می‌گیرد،
\dmterm{\textbf{عامل پرت محلی}}{Local Outlier Factor (LOF)} نام دارد.
مزیت اصلی LOF نسبت به روش‌های ساده‌تر این است که چگالی هر نقطه را نسبت
به چگالی همسایگانش می‌سنجد، نه نسبت به یک آستانهٔ سراسری ثابت؛ این
ویژگی به آن اجازه می‌دهد داده‌های پرت را حتی در مجموعه‌هایی که چگالی
خوشه‌های مختلف با هم تفاوت دارد (که در آن‌ها روش‌های سراسری‌تر دچار خطا
می‌شوند)، به‌درستی شناسایی کند.

\section{روش‌های مبتنی بر یادگیری}

\dmterm{\textbf{جنگل ایزوله}}{Isolation Forest} یکی از مؤثرترین
روش‌های مدرن تشخیص داده پرت است. ایدهٔ اصلی آن برخلاف روش‌های قبلی که
بر پایهٔ فاصله یا چگالی «عادی بودن» را می‌سنجیدند، مستقیماً بر پایهٔ
سهولت \textbf{ایزوله‌کردن} یک نمونه است: با ساخت درخت‌های تصادفی که
به‌طور مکرر داده را بر اساس ویژگی‌ها و مقادیر تصادفی تقسیم می‌کنند،
نمونه‌های پرت — که از بقیهٔ داده متمایزند — معمولاً با تعداد تقسیم بسیار
کمتری ایزوله می‌شوند (یعنی در عمق کمتری از درخت به تنهایی قرار
می‌گیرند)، در حالی که نمونه‌های عادی برای ایزوله شدن به تقسیم‌های بیشتری
نیاز دارند. میانگین عمق ایزوله‌سازی روی چند درخت تصادفی، امتیاز پرت
بودن هر نمونه را می‌دهد.

مزیت اصلی جنگل ایزوله نسبت به روش‌های مبتنی بر فاصله (مانند LOF)، کارایی
محاسباتی بهتر آن روی داده‌های با ابعاد بالا و حجم زیاد است، زیرا نیازی
به محاسبهٔ فاصلهٔ جفت‌به‌جفت میان همهٔ نمونه‌ها ندارد.

\section*{پیاده‌سازی در پایتون}

پیاده‌سازی \texttt{IsolationForest} و \texttt{LocalOutlierFactor} از
\texttt{scikit-learn}، به‌همراه محاسبهٔ امتیاز $Z$ و دامنهٔ میان‌چارکی
با \texttt{numpy}/\texttt{pandas}، در بخش «فصل ۸» پیوست الف
(\ref{app:python}) آمده است.

\section*{خلاصهٔ فصل}

در این فصل با مفهوم داده پرت، سه نوع اصلی آن (سراسری، زمینه‌ای، جمعی)
و طیفی از روش‌های تشخیص آن — از روش‌های آماری ساده (امتیاز $Z$، دامنهٔ
میان‌چارکی) تا روش‌های مبتنی بر چگالی (\LR{LOF}) و روش‌های مدرن مبتنی
بر یادگیری (جنگل ایزوله) — آشنا شدیم. با این فصل، بخش «یادگیری بدون
نظارت» این کتاب تکمیل می‌شود. در فصل بعد، به‌طور مختصر به دو حوزهٔ
تکمیلی، یعنی داده‌کاوی متن و وب، می‌پردازیم.

\section*{تمرین‌ها}

\begin{enumerate}
    \item برای هر یک از مثال‌های زیر مشخص کنید که به کدام نوع داده پرت
          (سراسری، زمینه‌ای یا جمعی) نزدیک‌تر است: (الف) خرید ناگهانی
          و بسیار بزرگ روی یک کارت اعتباری در نیمهٔ شب، (ب) یک روز با
          دمای ۳۰ درجه در وسط زمستان، (ج) یک سری کوتاه از درخواست‌های
          شبکه با الگوی غیرعادی که نشانهٔ حمله است.
    \item چرا روش‌های آماری تک‌متغیره (مانند امتیاز $Z$) برای داده‌های
          چندبعدی با روابط پیچیده میان ویژگی‌ها کافی نیستند؟
    \item تفاوت اصلی \LR{LOF} با روش‌های ساده‌تر مبتنی بر فاصله (مانند
          فاصله تا $k$اُمین همسایهٔ نزدیک) را در مواجهه با خوشه‌هایی با
          چگالی متفاوت توضیح دهید.
    \item ایدهٔ اصلی جنگل ایزوله را در یک یا دو جمله توضیح دهید و بگویید
          چرا این ایده برای داده‌های با ابعاد بالا کارآمد است.
    \item یک سناریوی واقعی (غیر از مثال‌های این فصل) پیدا کنید که در آن
          تشخیص داده پرت اهمیت عملی داشته باشد و توضیح دهید کدام‌یک از
          روش‌های این فصل برای آن مناسب‌تر است.
\end{enumerate}
